\begin{table}[H]\centering
\caption{\textbf{Municipality-year summary statistics}}\label{tab:descriptives}
\small
\small\setlength{\tabcolsep}{4pt}\resizebox{\ifdim\width>\linewidth \linewidth\else\width\fi}{!}{%
\begin{tabular}{lrrrrrr}
\toprule
Variable & Mean & Std.\ dev. & 10th pct. & Median & 90th pct. & $N$ \\
\midrule
Cesarean rate, private (TISS) & 0.84 & 0.09 & 0.73 & 0.85 & 0.94 & 5,198 \\
Cesarean rate, for-profit (SINASC) & 0.82 & 0.11 & 0.68 & 0.83 & 0.95 & 2,823 \\
Log economic fee gap (cesarean minus vaginal) & -0.04 & 0.38 & -0.42 & -0.08 & 0.44 & 5,198 \\
Obstetricians per 1,000 births & 24.02 & 70.92 & 11.09 & 20.97 & 35.09 & 5,189 \\
GDP per capita (R\$ thousands) & 65.97 & 30.37 & 33.71 & 62.13 & 91.88 & 4,715 \\
Private health-plan coverage (\%) & 38.78 & 12.25 & 21.23 & 40.02 & 53.25 & 5,198 \\
Adequate prenatal care (\%) & 70.31 & 10.66 & 55.10 & 73.12 & 81.39 & 5,198 \\
\bottomrule
\end{tabular}
}
\begin{minipage}{\linewidth}\footnotesize
\textit{Notes:} Municipality-years 2015--2024 with at least 20 private deliveries and a defined fee gap, the sample of the fee regressions, weighted by private deliveries as the regressions are; percentiles are weighted. TISS variables from ANS private-insurance claims; SINASC and covariates as described in the text. Unweighted, the mean log fee gap is +0.21: the gap is positive in most small municipalities, while 57 percent of deliveries take place where the cesarean pays less.
\end{minipage}
\end{table}
